%% ============================================================================
%% Campus Land Cover Classification with Sentinel-2 — White Paper
%% ============================================================================
\documentclass[12pt, a4paper, twocolumn]{article}

%% ── Packages ────────────────────────────────────────────────────────────────
\usepackage[utf8]{inputenc}
\usepackage[T1]{fontenc}
\usepackage{lmodern}
\usepackage[margin=2cm]{geometry}
\usepackage{graphicx}
\usepackage{amsmath, amssymb}
\usepackage{booktabs}
\usepackage{multirow}
\usepackage{caption}
\usepackage{subcaption}
\usepackage{hyperref}
\usepackage{xcolor}
\usepackage{float}
\usepackage{enumitem}
\usepackage{algorithm}
\usepackage{algpseudocode}
\usepackage{fancyhdr}
\usepackage{titlesec}
\usepackage{setspace}
\usepackage{url}
\usepackage{cite}
\usepackage{tabularx}
\usepackage{array}

%% ── Page Style ──────────────────────────────────────────────────────────────
\pagestyle{fancy}
\fancyhf{}
\fancyhead[L]{\small\textit{Campus Land Cover Classification with Sentinel-2}}
\fancyhead[R]{\small\thepage}
\renewcommand{\headrulewidth}{0.4pt}

%% ── Hyperref Colours ────────────────────────────────────────────────────────
\hypersetup{
    colorlinks=true,
    linkcolor=blue!70!black,
    citecolor=green!50!black,
    urlcolor=blue!60!black
}

%% ── Title ───────────────────────────────────────────────────────────────────
\title{
    \vspace{-1cm}
    \textbf{\Large Land Cover Classification of Campus and District-Scale Regions Using Sentinel-2 Imagery and Multi-Model Machine Learning on Google Earth Engine}
    \vspace{0.3cm}
}

\author{
    Ashutosh Saxena \\ \small\texttt{24bcs041@iiitdmj.ac.in} \and 
    Aarav Jain \\ \small\texttt{24bcs004@iiitdmj.ac.in} \and 
    Aaditya Poddar \\ \small\texttt{24bcs002@iiitdmj.ac.in} \and 
    Abhinav Kumar \\ \small\texttt{24bcs007@iiitdmj.ac.in} \and 
    Siddharth Ilayaraja \\ \small\texttt{24bcs292@iiitdmj.ac.in}
}

\date{June 2026}

%% ============================================================================
\begin{document}
\onecolumn
\maketitle
\begin{abstract}
\noindent
Accurate and scalable land cover classification is essential for environmental monitoring, urban planning, and natural resource management. This paper presents a comprehensive two-phase study on land cover mapping over the IIITDM Jabalpur university campus using multispectral satellite imagery from the Sentinel-2 mission. In \textbf{Phase~1} (binary classification), five supervised machine learning classifiers---Random Forest (RF), Support Vector Machine (SVM-RBF), Gradient Boosted Trees (XGBoost), K-Nearest Neighbours (KNN), and Classification and Regression Tree (CART)---were evaluated for \textit{Forest} vs.\ \textit{Non-Forest} discrimination using spectral bands B4, B8, and NDVI. Cloud masking was performed via QA60 bitmask processing and median compositing. Hyperparameter tuning used exhaustive grid search with 5-fold cross-validation. External validation on a spatially disjoint test region (480 points, 240 per class) showed XGBoost achieving the highest accuracy of 87.77\% ($\kappa = 0.75$), followed by SVM at 84.79\% ($\kappa = 0.70$) and RF at 81.87\% ($\kappa = 0.64$). District-scale deployment over Jabalpur ($\sim$5,211\,km$^2$, 1,600 auto-generated points) yielded accuracies exceeding 97\% for all models, with SVM attaining the highest at 98.31\% ($\kappa = 0.97$). In \textbf{Phase~2} (multi-class), the pipeline was extended to four semantic classes (\textit{Water}, \textit{Forest}, \textit{Soil}, \textit{Buildings}) using an expanded 10-dimensional feature space (B2, B3, B4, B8, B11, B12, NDVI, NDWI, NDBI, SAVI). All five classifiers were benchmarked via 5-fold cross-validation; XGBoost achieved the highest overall accuracy of 93.06\% ($\kappa = 0.892$), followed by SVM at 90.74\% ($\kappa = 0.873$), RF at 89.81\% ($\kappa = 0.859$), CART at 88.89\% ($\kappa = 0.847$), and KNN at 88.89\% ($\kappa = 0.847$). All models were executed server-side on Google Earth Engine with decoupled testing logic for full reproducibility.

\vspace{0.3cm}
\noindent\textbf{Keywords:} Land Cover Classification, Sentinel-2, Random Forest, SVM, XGBoost, KNN, CART, Google Earth Engine, NDVI, NDWI, NDBI, Multi-Class, Binary Classification
\end{abstract}

\twocolumn

%% ============================================================================
\section{Introduction}
\label{sec:introduction}
%% ============================================================================

Land cover classification is a fundamental task in remote sensing and geospatial analysis, underpinning applications in environmental monitoring, urban expansion tracking, forest management, and biodiversity assessment~\cite{lu2007survey}. Traditional manual surveys are expensive, time-consuming, and inherently non-scalable, motivating the development of automated machine-learning pipelines that can process large volumes of satellite imagery with minimal human intervention.

The advent of freely available, high-resolution satellite data---particularly from the European Space Agency's Sentinel-2 mission~\cite{drusch2012sentinel}---has democratised access to multispectral imagery at 10-metre spatial resolution with a five-day revisit cycle. Coupled with cloud-native geospatial computing platforms such as Google Earth Engine (GEE)~\cite{gorelick2017google}, researchers can now perform large-scale classification entirely server-side, eliminating the need for local data downloads and substantial computational infrastructure.

This study addresses two classification problems over the IIITDM Jabalpur campus (centred at approximately $80.02^{\circ}$E, $23.17^{\circ}$N): binary \textit{Forest} vs.\ \textit{Non-Forest} delineation, and four-class semantic land cover mapping (\textit{Water}, \textit{Forest}, \textit{Soil}, \textit{Buildings}). The primary objectives are:

\begin{enumerate}[label=(\roman*)]
    \item To benchmark five supervised classifiers---Random Forest, SVM (RBF kernel), XGBoost, KNN, and CART---for binary and multi-class land cover mapping using Sentinel-2 Surface Reflectance data.
    \item To assess model generalization by validating on an entirely unseen, spatially disjoint external test region using 480 independently collected test points.
    \item To evaluate district-scale spatial transferability by deploying trained models over the entire Jabalpur district boundary using 1,600 auto-generated ground-truth points.
    \item To perform systematic hyperparameter optimization via grid search with 5-fold cross-validation.
    \item To quantify the impact of seasonal phenology and cloud cover thresholds on classification accuracy across both binary and multi-class pipelines.
    \item To benchmark all five classifiers across per-class accuracy metrics (Forest, Water, Soil, Buildings) and summarise model-level cross-validation performance.
\end{enumerate}

The remainder of this paper is organized as follows: Section~\ref{sec:study_area} describes the study area. Section~\ref{sec:data} details data acquisition and preprocessing. Section~\ref{sec:methodology} presents the five classification algorithms. Section~\ref{sec:results} reports binary and district-level experimental results. Section~\ref{sec:multiclass} covers the multi-class extension and 5-model benchmark, and Section~\ref{sec:conclusion} concludes with key findings and future directions.

%% ============================================================================
\section{Study Area}
\label{sec:study_area}
%% ============================================================================

The study area is a university campus located in central India, defined by a five-vertex polygon boundary with the following coordinates:

\begin{itemize}
    \item $(80.0171^{\circ}\text{E},\; 23.1740^{\circ}\text{N})$
    \item $(80.0326^{\circ}\text{E},\; 23.1654^{\circ}\text{N})$
    \item $(80.0365^{\circ}\text{E},\; 23.1725^{\circ}\text{N})$
    \item $(80.0268^{\circ}\text{E},\; 23.1817^{\circ}\text{N})$
    \item $(80.0154^{\circ}\text{E},\; 23.1770^{\circ}\text{N})$
\end{itemize}

The campus encompasses a heterogeneous landscape comprising dense deciduous forest patches, built-up areas, open fields, and water bodies, providing a suitably complex environment for evaluating binary classification models.

An external validation region (centred at approximately $79.88^{\circ}$E, $23.14^{\circ}$N) was selected entirely outside the campus boundary. This spatially disjoint test region ensures that externally reported accuracies reflect true generalization capability rather than spatial autocorrelation effects.

\subsection{Jabalpur District Extension}
\label{subsec:jabalpur}

To evaluate district-scale transferability, the trained models were additionally deployed over the entire Jabalpur district (centred at approximately $79.98^{\circ}$E, $23.23^{\circ}$N), extracted from the \texttt{WM/geoLab/geoBoundaries/600/ADM2} FeatureCollection filtered by \texttt{shapeGroup = `IND'} and \texttt{shapeName = `Jabalpur'}. Jabalpur district spans approximately 5,211\,km$^2$ and presents a diverse mosaic of dense tropical dry deciduous forests (including portions of the Narmada river basin vegetation corridors), agricultural plains, built-up urban zones, and water bodies. This district-level deployment represents a significantly larger and more heterogeneous landscape compared to the campus-scale training area, providing a rigorous test of spatial transferability.

%% ============================================================================
\section{Data Acquisition and Preprocessing}
\label{sec:data}
%% ============================================================================

\subsection{Satellite Imagery}

The study utilises Sentinel-2 Level-2A Surface Reflectance (SR) imagery from the \texttt{COPERNICUS/S2\_SR\_HARMONIZED} collection available on Google Earth Engine. For the binary phase (Phase~1), two reflectance bands and one derived index were used; the multi-class phase (Phase~2) expanded the feature space to ten dimensions (see Section~\ref{sec:multiclass}). The binary phase features are:

\begin{table}[H]
\centering
\caption{Spectral features used for classification.}
\label{tab:bands}
\resizebox{\columnwidth}{!}{%
\begin{tabular}{@{}lll@{}}
\toprule
\textbf{Feature} & \textbf{Wavelength} & \textbf{Description} \\
\midrule
B4 (Red) & 665\,nm & Visible red reflectance \\
B8 (NIR) & 842\,nm & Near-infrared reflectance \\
NDVI & Derived & $\frac{\text{B8} - \text{B4}}{\text{B8} + \text{B4}}$ \\
\bottomrule
\end{tabular}%
}
\end{table}

The Normalized Difference Vegetation Index (NDVI)~\cite{rouse1974monitoring} serves as the primary discriminant feature, exploiting the strong Near-Infrared reflectance of healthy vegetation canopies against the Red absorption by chlorophyll pigments.

\subsection{Cloud Masking}

A two-stage cloud filtering strategy was employed (implemented in the \texttt{mask\_s2\_clouds} function):

\begin{enumerate}[label=\arabic*.]
    \item \textbf{Metadata-level filtering:} Only scenes below a specified \texttt{CLOUDY\_PIXEL\_PERCENTAGE} threshold were retained. The architecture strictly demands a maximum threshold of $\leq 20\%$.
    
    \item \textbf{Pixel-level masking:} Each retained image was processed using the QA60 bitmask band. Bits~10 (opaque clouds) and 11~(cirrus) were used to generate a binary mask:
    \begin{equation}
        M(x,y) = \overline{Q_{10}(x,y)} \wedge \overline{Q_{11}(x,y)}
    \end{equation}
    where $Q_{10}$ and $Q_{11}$ denote the respective QA60 bit values at pixel $(x,y)$. Masked pixels were excluded from the subsequent median composite.
\end{enumerate}

After masking, all pixel values were scaled by a factor of $1/10{,}000$ to convert from digital numbers to physical reflectance units.

\subsection{Image Compositing}

A temporal median composite was generated from all filtered images within the selected date range and clipped to the Region of Interest (ROI). Based on our rigorous phenological baseline discoveries, all models were trained against the dry season:

\begin{table}[H]
\centering
\caption{Date ranges and cloud thresholds per classifier notebook.}
\label{tab:dates}
\resizebox{\columnwidth}{!}{%
\begin{tabular}{@{}llll@{}}
\toprule
\textbf{Notebook} & \textbf{Date Range} & \textbf{Cloud \%} & \textbf{Seed} \\
\midrule
\texttt{nb/rf/} & 2026-01-01 to 2026-02-28 & 20\% & 42 \\
\texttt{nb/svm/} & 2026-01-01 to 2026-02-28 & 20\% & 120 \\
\texttt{nb/xgb/} & 2026-01-01 to 2026-02-28 & 20\% & 10 \\
\bottomrule
\end{tabular}%
}
\end{table}

For the Jabalpur district deployment, the composite imagery was generated over the date range \textbf{November 01 -- December 31, 2025} with a 20\% cloud pixel threshold, clipped to the district administrative boundary.

\subsection{Training Data}

Ground-truth training data consisted of 500 manually labelled points stored as GEE FeatureCollections:

\begin{itemize}
    \item \textbf{Forest:} 250 points (\texttt{label = 1}) via \url{users/cosypix/forest_points}
    \item \textbf{Non-Forest:} 250 points (\texttt{label = 0}) via \url{users/cosypix/non_forest_points}
\end{itemize}

Point counts were evaluated and passed directly into the sampling architectures spanning a 70/30 random split utilizing \texttt{training.randomColumn(`random', SEED)}. Null-value samples were excluded via \texttt{ee.Filter.notNull}.

For external validation, an independent set of 480 labelled points (240 per class) was rigorously deployed within the spatially disjoint test region:

\begin{itemize}
    \item \textbf{Test Forest:} 240 points (\texttt{label = 1}) via \url{users/cosypix/test_forest_points}
    \item \textbf{Test Non-Forest:} 240 points (\texttt{label = 0}) via \url{users/cosypix/test_non_forest_points}
\end{itemize}

\subsection{Jabalpur District Test Data}
\label{subsec:jbp_data}

For the Jabalpur district-level validation, a separate set of 1,600 auto-generated test points (800 per class) was deployed across the full district extent:

\begin{itemize}
    \item \textbf{Test Forest:} 800 points (\texttt{label = 1}) via \url{users/cosypix/jabalpurforestnew}
    \item \textbf{Test Non-Forest:} 800 points (\texttt{label = 0}) via \url{users/cosypix/jabalpurnonforestnew}
\end{itemize}

These test points were auto-generated using high-confidence reference labels derived from ESA WorldCover / Dynamic World products, stratified-randomly distributed across the district boundary, and exported as GEE assets. The substantially larger test set (1,600 vs.\ 480 points) provides enhanced statistical power for evaluating model performance at the district scale.

%% ============================================================================
\section{Methodology}
\label{sec:methodology}
%% ============================================================================

Five supervised classification algorithms were selected to represent fundamentally different learning paradigms: bagging (Random Forest), margin-based learning (SVM), boosting (XGBoost), instance-based learning (KNN), and single decision tree (CART). For the binary phase, classifiers were trained on the three-band feature vector $\mathbf{x} = [\text{B4}, \text{B8}, \text{NDVI}]^\top$; for the multi-class phase the feature space was expanded to ten dimensions (see Section~\ref{sec:multiclass}).

\subsection{Random Forest (RF)}
\label{subsec:rf}

Random Forest~\cite{breiman2001random} constructs multiple decision trees during training, each grown on a bootstrapped sample. \textbf{GEE Implementation:} The \texttt{ee.Classifier.smileRandomForest} classifier was used. The canonical notebook uses:

\begin{table}[H]
\centering
\caption{Random Forest -- Optimized Hyperparameters.}
\label{tab:rf_params}
\begin{tabular}{@{}ll@{}}
\toprule
\textbf{Parameter} & \textbf{Value} \\
\midrule
\texttt{numberOfTrees} & 10 \\
\texttt{minLeafPopulation} & 3 \\
\texttt{bagFraction} & 0.3 \\
\bottomrule
\end{tabular}
\end{table}

\subsection{Support Vector Machine (SVM-RBF)}
\label{subsec:svm}

Support Vector Machines~\cite{cortes1995support} seek the maximum-margin hyperplane that separates classes. The \textbf{GEE Implementation:} \texttt{ee.Classifier.libsvm} classifier uses:

\begin{table}[H]
\centering
\caption{SVM -- Optimized Hyperparameters.}
\label{tab:svm_params}
\begin{tabular}{@{}ll@{}}
\toprule
\textbf{Parameter} & \textbf{Value} \\
\midrule
\texttt{kernelType} & RBF \\
\texttt{gamma} ($\gamma$) & 0.5 \\
\texttt{cost} ($C$) & 10.0 \\
\bottomrule
\end{tabular}
\end{table}

\subsection{Gradient Boosted Trees (XGBoost)}
\label{subsec:xgb}

Gradient Boosted Trees~\cite{chen2016xgboost} sequentially construct an ensemble of weak learners (shallow decision trees). The \textbf{GEE Implementation:} \texttt{ee.Classifier.smileGradientTreeBoost} classifier utilizes:

\begin{table}[H]
\centering
\caption{XGBoost -- Optimized Hyperparameters.}
\label{tab:xgb_params}
\begin{tabular}{@{}ll@{}}
\toprule
\textbf{Parameter} & \textbf{Value} \\
\midrule
\texttt{numberOfTrees} & 10 \\
\texttt{shrinkage} ($\eta$) & 0.01 \\
\texttt{maxNodes} & 3 \\
\bottomrule
\end{tabular}
\end{table}

\subsection{K-Nearest Neighbours (KNN)}
\label{subsec:knn}

K-Nearest Neighbours is a non-parametric, instance-based classifier that assigns a label by majority vote among the $k$ closest training samples in feature space. \textbf{GEE Implementation:} \texttt{ee.Classifier.minimumDistance} with the Euclidean metric was used. For the multi-class phase, the optimal neighbourhood size was determined via grid search cross-validation.

\begin{table}[H]
\centering
\caption{KNN -- Optimized Hyperparameters (Multi-Class).}
\label{tab:knn_params}
\begin{tabular}{@{}ll@{}}
\toprule
\textbf{Parameter} & \textbf{Value} \\
\midrule
\texttt{metric} & Euclidean \\
\texttt{k} & tuned via CV \\
\bottomrule
\end{tabular}
\end{table}

\subsection{Classification and Regression Tree (CART)}
\label{subsec:cart}

CART is a single decision tree algorithm that recursively partitions the feature space using binary splits to minimise impurity (Gini index for classification). \textbf{GEE Implementation:} \texttt{ee.Classifier.smileCart} was used. CART provides a human-interpretable model and acts as a structural baseline against ensemble methods.

\begin{table}[H]
\centering
\caption{CART -- Optimized Hyperparameters (Multi-Class).}
\label{tab:cart_params}
\begin{tabular}{@{}ll@{}}
\toprule
\textbf{Parameter} & \textbf{Value} \\
\midrule
\texttt{maxNodes} & tuned via CV \\
\texttt{minLeafPopulation} & tuned via CV \\
\bottomrule
\end{tabular}
\end{table}

%% ============================================================================
\section{Experimental Results}
\label{sec:results}
%% ============================================================================

\subsection{External Validation (Unseen Test Area)}

External validation was conducted on a spatially disjoint test region having 2 times the area of our college using 480 independently collected ground-truth points (240 per class). The cal\_area notebooks executed the models via global namespace \texttt{\%run} variables mapping to the tuned hyperparameter pipelines.

\begin{table}[H]
\centering
\caption{External validation results on unseen test area (240 per class).}
\label{tab:external}
\resizebox{\columnwidth}{!}{%
\begin{tabular}{@{}llcc@{}}
\toprule
\textbf{Classifier} & \textbf{Confusion Matrix} & \textbf{Acc.\ (\%)} & \textbf{$\kappa$} \\
\midrule
XGBoost       & {[[218,22],[22,218]]} & \textbf{87.77} & \textbf{0.7541} \\
SVM (RBF)     & {[[219,21],[22,218]]} & 84.79          & 0.6958 \\
Random Forest & {[[214,26],[26,214]]} & 81.87          & 0.6375 \\
\bottomrule
\end{tabular}%
}
\end{table}

The external test results demonstrate that \textbf{XGBoost} achieves the best performance on completely unseen data (87.77\%, $\kappa=0.754$), followed by \textbf{SVM (RBF)} (84.79\%, $\kappa=0.696$) and \textbf{Random Forest} (81.87\%, $\kappa=0.638$).

\subsection{District-Level Validation: Jabalpur}
\label{subsec:jbp_results}

To evaluate the spatial transferability of the campus-trained classifiers at a significantly larger administrative scale, all three binary-phase classifiers (RF, SVM, XGBoost) were deployed over the entire Jabalpur district boundary. A Sentinel-2 median composite (November--December 2025, $\leq 20\%$ cloud cover) was generated and clipped to the district geometry. Classification was performed using the pre-trained models loaded via \texttt{\%run} from the respective tuned notebooks. Post-classification spatial smoothing was applied using a \texttt{focalMode(1)} filter.

Validation was conducted against 1,600 independently auto-generated ground-truth points (800 Forest + 800 Non-Forest), stratified across the full district extent.

\begin{table}[H]
\centering
\caption{Jabalpur district-level validation results (800 per class, $N = 1{,}600$).}
\label{tab:jbp_results}
\resizebox{\columnwidth}{!}{%
\begin{tabular}{@{}llcc@{}}
\toprule
\textbf{Classifier} & \textbf{Confusion Matrix} & \textbf{Acc.\ (\%)} & \textbf{$\kappa$} \\
\midrule
SVM (RBF) & {[[785,15],[12,788]]} & \textbf{98.31} & \textbf{0.9663} \\
XGBoost & {[[794,6],[25,775]]} & 98.06 & 0.9613 \\
Random Forest & {[[789,11],[25,775]]} & 97.75 & 0.9550 \\
\bottomrule
\end{tabular}%
}
\end{table}

All three binary classifiers achieved district-level accuracies exceeding 97\%, substantially surpassing the campus-adjacent external validation scores. \textbf{SVM (RBF)} yielded the highest overall accuracy of 98.31\% ($\kappa = 0.9663$), followed by \textbf{XGBoost} at 98.06\% ($\kappa = 0.9613$) and \textbf{Random Forest} at 97.75\% ($\kappa = 0.9550$).

\subsubsection{Confusion Matrix Analysis -- Jabalpur}

Examining the per-class performance for the SVM classifier on the Jabalpur district:

\begin{itemize}
    \item \textbf{True Non-Forest rate (Specificity):} $785 / (785 + 15) = 98.13\%$
    \item \textbf{True Forest rate (Sensitivity):} $788 / (788 + 12) = 98.50\%$
    \item \textbf{Forest Commission Errors:} Only 15 Non-Forest pixels misclassified as Forest (1.88\%)
    \item \textbf{Forest Omission Errors:} Only 12 Forest pixels misclassified as Non-Forest (1.50\%)
\end{itemize}

The near-symmetric confusion matrices across all classifiers indicate minimal systematic bias and robust generalization across the heterogeneous Jabalpur landscape.

\subsubsection{Comparative Analysis: External vs.\ District-Level}

\begin{table}[H]
\centering
\caption{Accuracy comparison across validation scales.}
\label{tab:scale_comparison}
\resizebox{\columnwidth}{!}{%
\begin{tabular}{@{}lccc@{}}
\toprule
\textbf{Classifier} & \textbf{Campus Ext.\ (\%)} & \textbf{Jabalpur (\%)} & \textbf{$\Delta$ (\%)} \\
\midrule
SVM (RBF) & 84.79 & \textbf{98.31} & $+13.52$ \\
XGBoost & 87.77 & 98.06 & $+10.29$ \\
Random Forest & 81.87 & 97.75 & $+15.88$ \\
\bottomrule
\end{tabular}%
}
\end{table}

The pronounced accuracy improvements at district scale ($+10$ to $+16$ percentage points) are attributed to several factors: (i) the auto-generated Jabalpur test points leverage high-confidence reference labels from ESA WorldCover/Dynamic World, reducing labelling noise; (ii) the test point distribution spans a broader spatial extent, averaging out localised spectral anomalies; and (iii) the Jabalpur imagery date range (Nov--Dec 2025) corresponds to the post-monsoon dry transition, where forest canopy retention provides strong NDVI contrast against desiccated non-forest surfaces.

\subsection{Sensitivity Analysis \& Seasonal Stability}

An extensive temporal ablation study was performed to quantify model reliance on specific seasonal metadata. To identify the optimal phenomenological window for binary forest delineation, three distinct intervals were passed through the prediction pipeline targeting the testing area:

\begin{itemize}
    \item \textbf{Early Summer:} March 01 -- May 31, 2025
    \item \textbf{Post-Monsoon:} October 01 -- November 30, 2025
    \item \textbf{Winter (Baseline):} January 01 -- February 28, 2026
\end{itemize}

By restricting the median composite generation strictly to these discrete intervals, the isolated effect of seasonal canopy density and undergrowth moisture distribution was thoroughly quantified.

\begin{table}[H]
\centering
\caption{Seasonal accuracy comparisons demonstrating catastrophic classification collapses outside of the dry Winter season bounds.}
\label{tab:seasons}
\resizebox{\columnwidth}{!}{%
\begin{tabular}{@{}lccc@{}}
\toprule
\textbf{Classifier} & \textbf{Early Summer} & \textbf{Post-Monsoon} & \textbf{Winter} \\
\midrule
Standard Model & 88.13\% & 70.42\% & \textbf{90.00\%} \\
\bottomrule
\end{tabular}%
}
\end{table}

\textbf{Winter imagery consistently yielded the highest absolute stability scaling.} Internal diagnostics revealed that during the Post-Monsoon (Oct-Nov), massive distribution shifts occurred within the spectral layout. Wild grasses and agricultural weed plots spiked significantly in intensity, causing non-forest terrain values to effectively hijack the decision boundaries. During Jan-Feb, short-lived green vegetation dramatically dies off across the Central Indian topography, leaving permanently rooted forest trees as practically the lone surviving source of distinctly vibrant green space, radically improving classification contours.

Furthermore, strictly bounding the \texttt{CLOUDY\_PIXEL\_PERCENTAGE} filtering cutoff parameters to a tolerance of 20\% shielded the remaining classifiers from severe misclassification noise.

\subsection{Confusion Matrix Analysis}

Examining the SVM confusion matrix in detail on the external region:
\begin{itemize}
    \item \textbf{True Non-Forest rate (Specificity):} $219 / (219 + 21) = 91.25\%$
    \item \textbf{True Forest rate (Sensitivity):} $218 / (218 + 22) = 90.83\%$
    \item The matrix is highly symmetric, minimizing systematic label bias during classification tasks spanning external regions.
\end{itemize}

%% ============================================================================
\section{Multi-Class Extension}
\label{sec:multiclass}
%% ============================================================================

To provide a richer semantic understanding of the campus landscape, the binary classification pipeline was extended to a four-class scheme: \textit{Water} (0), \textit{Forest} (1), \textit{Soil/Open Ground} (2), and \textit{Buildings/Built-up} (3). All five classifiers (RF, SVM, XGBoost, KNN, CART) were trained and evaluated in this phase.

\subsection{Expanded Feature Space and Training Data}
Distinguishing four heterogeneous classes required expanding the spectral feature space from three to ten dimensions. The new feature vector $\mathbf{x}$ included: B2 (Blue), B3 (Green), B4 (Red), B8 (NIR), B11 (SWIR-1), B12 (SWIR-2), NDVI, Normalized Difference Water Index (NDWI), Normalized Difference Built-up Index (NDBI), and Soil Adjusted Vegetation Index (SAVI).

Ground-truth points were independently generated using Google Dynamic World and ESA WorldCover reference masks, yielding a robust dataset of approximately 130--150 points per class distributed across the campus geometry. The optimal imagery configuration---\textbf{Early Summer (March--May), 15\% cloud cover}---was identified through sensitivity analysis as it provides the strongest multi-class spectral separation.

\subsection{5-Model Cross-Validation Results}

All five classifiers were evaluated via 5-fold stratified cross-validation with grid search hyperparameter tuning. Table~\ref{tab:multi_overall} reports the overall accuracy and Cohen's kappa for each model, demonstrating the superiority of ensemble and kernel methods over single-tree and instance-based approaches.

\begin{table}[H]
\centering
\caption{Multi-class 5-fold CV results -- all 5 classifiers (campus dataset).}
\label{tab:multi_overall}
\resizebox{\columnwidth}{!}{%
\begin{tabular}{@{}lcc@{}}
\toprule
\textbf{Classifier} & \textbf{Accuracy (\%)} & \textbf{$\kappa$} \\
\midrule
XGBoost             & \textbf{93.06}         & \textbf{0.892} \\
SVM (RBF)           & 90.74                  & 0.873 \\
Random Forest       & 89.81                  & 0.859 \\
CART                & 88.89                  & 0.847 \\
KNN                 & 88.89                  & 0.847 \\
\bottomrule
\end{tabular}%
}
\end{table}

\textbf{XGBoost} achieved the best overall accuracy (93.06\%, $\kappa=0.892$) across all five classifiers. SVM (RBF) followed closely at 90.74\%, with RF at 89.81\%. CART and KNN tied at 88.89\%, confirming that ensemble methods outperform structural baselines on this heterogeneous four-class problem.

\subsection{Per-Class Geographical Accuracy}

Table~\ref{tab:multi_geo} provides the per-class producer's accuracy for each model across the four land cover categories, revealing distinct strengths and weaknesses by classifier type.

\begin{table}[H]
\centering
\caption{Per-class producer accuracy (\%) for all 5 classifiers.}
\label{tab:multi_geo}
\resizebox{\columnwidth}{!}{%
\begin{tabular}{@{}lcccc@{}}
\toprule
\textbf{Classifier} & \textbf{Forest} & \textbf{Water} & \textbf{Buildings} & \textbf{Soil} \\
\midrule
XGBoost       & \textbf{95.71} & \textbf{100.00} & 95.56 & 86.11 \\
SVM (RBF)     & 91.55          & 96.43           & 93.33 & 86.11 \\
Random Forest & 91.67          & \textbf{100.00} & 95.24 & 81.58 \\
KNN           & 87.36          & 93.33           & \textbf{96.00} & \textbf{90.24} \\
CART          & \textbf{100.00}& 94.20           & 84.06 & 83.64 \\
\bottomrule
\end{tabular}%
}
\end{table}

Key observations: (i) CART achieves perfect Forest recall (100\%) but is weakest on Buildings (84.06\%); (ii) XGBoost and RF achieve perfect Water recall (100\%), leveraging the strong NDWI signal; (iii) KNN demonstrates the highest Soil accuracy (90.24\%), benefiting from its local-distance sensitivity to SWIR-based Soil signatures; (iv) XGBoost provides the most balanced performance across all four classes.

\subsection{Cross-Validation Statistical Visualizations}

Figure~\ref{fig:cv_graphs} presents the 5-fold cross-validation accuracy trajectories, log-loss curves, and multi-class AUC-ROC curves for all five classifiers, providing a holistic statistical view of model stability across folds.

\begin{figure}[H]
    \centering
    \begin{subfigure}[b]{0.95\columnwidth}
        \includegraphics[width=\textwidth]{assets/images/cv_accuracy_comparison.png}
        \caption{5-Fold CV Accuracy per model and fold.}
    \end{subfigure}
    \vskip 0.3\baselineskip
    \begin{subfigure}[b]{0.95\columnwidth}
        \includegraphics[width=\textwidth]{assets/images/cv_loss_comparison.png}
        \caption{5-Fold CV Log-Loss per model and fold.}
    \end{subfigure}
    \vskip 0.3\baselineskip
    \begin{subfigure}[b]{0.95\columnwidth}
        \includegraphics[width=\textwidth]{assets/images/cv_roc_curve_comparison.png}
        \caption{Multi-class AUC-ROC curves (One-vs-Rest).}
    \end{subfigure}
    \caption{Cross-validation statistical performance for all 5 classifiers on the 4-class campus dataset.}
    \label{fig:cv_graphs}
\end{figure}

\subsection{Confusion Matrix Analysis -- All 5 Classifiers}

Figures~\ref{fig:multi_cm_ensemble} and~\ref{fig:multi_cm_baseline} illustrate the confusion matrices for all five classifiers on the multi-class campus validation set, demonstrating strong diagonal dominance across all categories.

\begin{figure}[H]
    \centering
    \begin{subfigure}[b]{0.45\columnwidth}
        \includegraphics[width=\textwidth]{assets/images/multi_rf_cm.png}
        \caption{Random Forest}
    \end{subfigure}
    \hfill
    \begin{subfigure}[b]{0.45\columnwidth}
        \includegraphics[width=\textwidth]{assets/images/multi_svm_cm.png}
        \caption{SVM (RBF)}
    \end{subfigure}
    \vskip\baselineskip
    \begin{subfigure}[b]{0.45\columnwidth}
        \includegraphics[width=\textwidth]{assets/images/multi_xgb_cm.png}
        \caption{XGBoost}
    \end{subfigure}
    \caption{Confusion matrices -- ensemble classifiers (RF, SVM, XGBoost).}
    \label{fig:multi_cm_ensemble}
\end{figure}

\begin{figure}[H]
    \centering
    \begin{subfigure}[b]{0.45\columnwidth}
        \includegraphics[width=\textwidth]{assets/images/multi_knn_cm.png}
        \caption{KNN}
    \end{subfigure}
    \hfill
    \begin{subfigure}[b]{0.45\columnwidth}
        \includegraphics[width=\textwidth]{assets/images/multi_cart_cm.png}
        \caption{CART}
    \end{subfigure}
    \caption{Confusion matrices -- baseline classifiers (KNN, CART).}
    \label{fig:multi_cm_baseline}
\end{figure}

\subsection{Multi-Class Sensitivity Analysis}

To ensure robustness, a combined sensitivity analysis was conducted across three distinct seasons (Early Summer, Post-Monsoon, and Winter) and six cloud cover thresholds (5\%--50\%). The \textbf{Early Summer (March--May) at 15\% cloud cover} configuration was confirmed as the optimal operating point, delivering the highest and most stable multi-class accuracies across all five models.

The expanded 10-dimensional feature space (incorporating SWIR bands and targeted indices NDWI/NDBI) significantly mitigated the seasonal degradation previously observed in the binary pipeline. While the binary model suffered classification collapse during Post-Monsoon due to agricultural weed interference, multi-class models maintained high stability. Specifically, SVM (RBF) proved the most resilient across seasons (Early Summer: 96.04\%, Post-Monsoon: 95.08\%, Winter: 95.33\% at 20\% cloud cover). XGBoost and SVM maintained $>$95\% accuracy even at 50\% cloud cover thresholds during Winter, confirming that the extended feature space substantially improves both seasonal generalizability and cloud cover tolerance.

%% ============================================================================
\section{Conclusion}
\label{sec:conclusion}

This study rigorously evaluated five supervised machine learning classifiers---Random Forest, SVM (RBF), XGBoost, KNN, and CART---across two phases of land cover classification using Sentinel-2 imagery on Google Earth Engine. The principal findings are:

\begin{enumerate}
    \item \textbf{Binary Model Generalization:} On an external spatially disjoint test region (480 points), XGBoost achieved the highest external accuracy of 87.77\% ($\kappa=0.75$), followed by SVM at 84.79\% and RF at 81.87\%. Nonlinear ensemble methods outperform structural baselines under spatial distribution shift.

    \item \textbf{District-Scale Transferability:} Deployed across the entire Jabalpur district ($\sim$5,211\,km$^2$, 1,600 auto-generated test points), all three primary classifiers exceeded 97\% accuracy. SVM (RBF) achieved the highest district accuracy of 98.31\% ($\kappa=0.97$), confirming that campus-trained models transfer effectively to district-scale landscapes within the same climatic zone.

    \item \textbf{Multi-Class Performance (5-Model Benchmark):} Across the full 4-class campus benchmark, XGBoost led with 93.06\% ($\kappa=0.892$), followed by SVM (90.74\%), RF (89.81\%), and CART and KNN (both 88.89\%). Ensemble methods consistently outperform instance-based and single-tree approaches on heterogeneous multi-class terrain.

    \item \textbf{Per-Class Strengths:} XGBoost and RF attain perfect Water recall (100\%) via strong NDWI contrast. CART achieves perfect Forest recall (100\%) but struggles on Buildings. KNN uniquely excels on Soil (90.24\%) due to local distance sensitivity in SWIR space.

    \item \textbf{Optimal Configuration:} Early Summer (March--May) imagery at 15\% cloud cover was identified as the optimal multi-class operating point. Winter (Jan--Feb) remains the best binary classification window due to maximum undergrowth senescence and forest-non-forest spectral contrast.

    \item \textbf{Seasonal \& Cloud Robustness:} The 10-dimensional multi-class feature space (NDWI, NDBI, SAVI, SWIR) substantially reduced the seasonal degradation observed in the binary pipeline, enabling $>$95\% accuracy across all seasons and cloud thresholds up to 50\%.
\end{enumerate}

\subsection{Future Work}

\begin{itemize}
    \item \textbf{Multi-sensor fusion:} Integrating Sentinel-1 SAR backscatter data to provide canopy structure cues independent of optical phenological conditions, directly addressing forest-vs-agricultural confusion in the Post-Monsoon window.
    \item \textbf{Multi-district scaling:} Extending the district-level validation framework to additional administrative regions across Madhya Pradesh to quantify cross-district transferability bounds.
    \item \textbf{Deep learning integration:} Incorporating convolutional neural networks (e.g., U-Net) for pixel-wise semantic segmentation to potentially surpass the 93\% multi-class ceiling achieved by XGBoost.
    \item \textbf{Temporal stacking:} Building multi-temporal feature stacks from all three seasonal composites simultaneously to exploit phenological trajectories as classification features.
\end{itemize}

%% ============================================================================
%% REFERENCES
%% ============================================================================
\begin{thebibliography}{99}

\bibitem{lu2007survey}
D.~Lu and Q.~Weng, ``A survey of image classification methods and techniques for improving classification performance,'' \textit{International Journal of Remote Sensing}, vol.~28, no.~5, pp.~823--870, 2007.

\bibitem{drusch2012sentinel}
M.~Drusch, U.~Del~Bello, S.~Carlier, \textit{et~al.}, ``Sentinel-2: ESA's optical high-resolution mission for GMES operational services,'' \textit{Remote Sensing of Environment}, vol.~120, pp.~25--36, 2012.

\bibitem{gorelick2017google}
N.~Gorelick, M.~Hancher, M.~Dixon, S.~Ilyushchenko, D.~Thau, and R.~Moore, ``Google Earth Engine: Planetary-scale geospatial analysis for everyone,'' \textit{Remote Sensing of Environment}, vol.~202, pp.~18--27, 2017.

\bibitem{rouse1974monitoring}
J.~W.~Rouse, R.~H.~Haas, J.~A.~Schell, and D.~W.~Deering, ``Monitoring vegetation systems in the Great Plains with ERTS,'' in \textit{Proc.\ Third Earth Resources Technology Satellite-1 Symposium}, pp.~309--317, 1974.

\bibitem{breiman2001random}
L.~Breiman, ``Random forests,'' \textit{Machine Learning}, vol.~45, no.~1, pp.~5--32, 2001.

\bibitem{cortes1995support}
C.~Cortes and V.~Vapnik, ``Support-vector networks,'' \textit{Machine Learning}, vol.~20, no.~3, pp.~273--297, 1995.

\bibitem{vapnik1998statistical}
V.~N.~Vapnik, \textit{Statistical Learning Theory}. New York: Wiley, 1998.

\bibitem{friedman2001greedy}
J.~H.~Friedman, ``Greedy function approximation: A gradient boosting machine,'' \textit{Annals of Statistics}, vol.~29, no.~5, pp.~1189--1232, 2001.

\bibitem{chen2016xgboost}
T.~Chen and C.~Guestrin, ``XGBoost: A scalable tree boosting system,'' in \textit{Proc.\ 22nd ACM SIGKDD International Conference on Knowledge Discovery and Data Mining}, pp.~785--794, 2016.

\bibitem{fix1951discriminatory}
E.~Fix and J.~L.~Hodges, ``Discriminatory analysis: Nonparametric discrimination: Consistency properties,'' \textit{USAF School of Aviation Medicine Technical Report}, vol.~4, pp.~261--279, 1951.

\bibitem{breiman1984classification}
L.~Breiman, J.~H.~Friedman, R.~A.~Olshen, and C.~J.~Stone, \textit{Classification and Regression Trees}. Belmont, CA: Wadsworth, 1984.

\end{thebibliography}

\end{document}

